\oldpage[IV-2]{181}
\begin{proposition}[6.15.10]
\label{IV.2.s6.15.10-fr}
Sean $Y$ y $Y'$ dos preesquemas localmente noetherianos,
$g:Y'\to Y$ un morfismo normal \textup{(6.8.1)}, y $f:X\to Y$ un
morfismo; pongamos $X'=X\times_Y Y'$ y sean $p:X'\to X$
la proyección canónica y $x'$ un punto de $X'$. Entonces,
si $X$ es reducido \textup{(resp.} geométricamente unirramoso,
\textup{resp.} íntegro y geométricamente unirramoso\textup{)} en el punto
$x=p(x')$, $X'$ es reducido \textup{(resp.} geométricamente
unirramoso, \textup{resp.} íntegro y geométricamente unirramoso\textup{)}
en el punto $x'$.
\end{proposition}

\begin{proof}
En virtud de (\textbf I, 3.6.5), es posible limitarse al caso en que
$Y=\operatorname{Spec}(A)$, $Y'=\operatorname{Spec}(A')$ y
$X=\operatorname{Spec}(B)$, donde $A$, $A'$ y $B$ son anillos
locales, los homomorfismos $A\to A'$ y $A\to B$ son locales,
$A$ y $A'$ son anillos noetherianos, y $x$ es el punto cerrado
de $X$. Se trata de probar que, si $B$ es reducido (resp.
geométricamente unirramoso, resp. íntegro y geométricamente
unirramoso), entonces $B'=B\otimes_A A'$ es reducido (resp.
$\operatorname{Spec}(B')$ es geométricamente unirramoso en
los puntos de $p^{-1}(x)$, resp. $\operatorname{Spec}(B')$
es íntegro y geométricamente unirramoso en esos puntos).
Supongamos primero únicamente que $B$ es reducido; $B$ es
la unión de la familia creciente de sus sub-$A$-álgebras
de tipo finito $B_\lambda$ y, por planitud, $B'$ es la unión
de la familia filtrante creciente de sus sub-$A'$-álgebras
$B'_\lambda=B_\lambda\otimes_A A'$. Ahora bien, el morfismo
$\operatorname{Spec}(B'_\lambda)\to\operatorname{Spec}(B_\lambda)$
es normal (6.8.2), y \emph{con mayor razón} reducido; el hecho de
que $B'_\lambda$ sea reducido se sigue, por tanto, de (3.3.5).
Se concluye por (5.13.2) que $B'$ mismo es reducido.

Supongamos ahora que $B$ es geométricamente unirramoso;
teniendo en cuenta (\textbf{I}, 5.1.8), se puede suponer además que
$B$ es reducido y, por tanto, \emph{íntegro}. Sea $C$ su clausura
integral, que por hipótesis es un anillo \emph{local}; pongamos
$Z=\operatorname{Spec}(C)$, $C'=C\otimes_A A'$ y
$Z'=\operatorname{Spec}(C')=Z\times_X X'$, de modo que se
tiene el diagrama conmutativo
\[
\xymatrix{
Z \ar[d]_f & Z' \ar[l] \ar[d]^{f'}\\
X & X' \ar[l]^p
}
\]
En virtud de la primera parte del razonamiento, $X'$ es
\emph{reducido}; por otra parte, como $Z'=Z\times_Y Y'$, se sigue
de (6.14.1) que $Z'$ es \emph{normal} (y \emph{con mayor razón},
geométricamente unirramoso). Como $l$ es entero y
birracional, lo mismo vale para $l'$ por (6.15.4.1), ya
que $p$ es plano. Finalmente, la hipótesis de que $X$ es
geométricamente unirramoso en $x$ implica que $l$ es
radicial en ese punto (6.15.5); se concluye, por tanto, de
(6.15.3.1) que $l'$ es radicial en todo punto
$x'\in p^{-1}(x)$, y de (6.15.5) se sigue que $X'$ es
geométricamente unirramoso (y, por tanto, íntegro, ya que
es reducido) en cada uno de esos puntos.
\end{proof}
